\chapter{Innvendingane}

\begin{motto}
Eg har lese det sterkaste forsvaret faget kan mønstre.\\ Her er det, i si beste form, og her er kvar det held meg.
\end{motto}

Ei bok som denne har ei plikt ho ikkje kan kjøpe seg fri frå: å møte den beste motstanden og ikkje den dårlegaste. Det finst eit seriøst forsvar for designfaget, ført av kloke folk over fleire tiår, og om eg berre felte stråmenn var heile klagemålet verdlaust. Så her reiser eg dei sterkaste innvendingane mot tesen i denne boka, kvar i si mest overtydande form, og prøver henne med éin einaste kniv. Kniven er ikkje «kvifor er ikkje design som fysikk?». Det spørsmålet er ein blindveg, og fleire av innvendingane har rett i å avvise det. Kniven er smalare: kvar er designens eige delte, prøvbare grunnlag for å seie kvifor \emph{denne} forma er betre enn \emph{den}? Eg stiller berre kniven her. Eg landar henne ikkje, for nokre av innvendingane bit tilbake, og dei skal få lov.

\section{«Kravet er ein kategorifeil»}

Den fyrste innvendinga er òg den som har mest rett, og ho fortener å bli ført i kjøt. Sjå føre deg ein torsdag, eit auditorium, regn på rutene, og den dyktigaste forsvararen faget eig reiser seg utan manus etter at eg har sett meg. «Du måler oss mot ein standard vi aldri lova å oppfylle,» seier han. Design høyrer til vitskapane om det kunstige, om korleis ting \emph{kan} bli og ikkje korleis dei \emph{er} \autocite{simon1969sciences}; det er ein eigen disiplin med eigne måtar å vite på, ikkje ein vitskap som kom til kort \autocite{cross1982, cross2001discipline}. Schön synte at den reflekterande praktikaren tenkjer \emph{i} handlinga, ikkje før henne, og at det er ein full kunnskapsform \autocite{schon1983}. «Kven sa at fysikk var den einaste forma for grunn? Du krev at vi skal sjå ut som noko vi ikkje er, og dømer oss så for ikkje å likne.»

Han har rett, og eg seier det utan vilkår. Kravet om at design skal grunngje forma slik fysikken grunngjev fakta, er ein kategorifeil, og eg skal ikkje smugle han inn att bakvegen ved å halde jus eller medisin fram som ein mildare mal design «burde» ha nått. Faget er ikkje ein mislukka vitskap. Det er eit anna slag ting. Men sjå kva som skjer når eg gjev innvendinga alt dette, for ho flyttar berre spørsmålet, ho slukkar det ikkje. At design ikkje skal vere fysikk, fortel meg ingenting om kvar design \emph{sjølv} finn grunnen til at éi form er rett. Forsvararen tala om Simon, om Cross, om Schön, og kvar av dei er ein teori om korleis ein designar \emph{tenkjer} eller \emph{arbeider}. Ingen av dei er ein delt prøve på den ferdige forma. Han førte fram alt rundt holet, vakkert, og bad salen sjå mengda og gløyme midten. Det er her eg ikkje kjem heilt i mål, og det skal stå: eg kan vise at innvendinga ikkje fyller holet, men eg kan ikkje vise at holet \emph{må} fyllast på det viset eg etterlyser. Kanskje er form det slag ting som ikkje toler eit slikt grunnlag. Då har forsvararen vunne meir enn han trur, og tapt på ein verre måte: faget han forsvarar, kviler då ikkje på ein grunn det enno manglar, men på ein grunn det aldri kan få. Köln i juli 1914 er det reine biletet. Muthesius la fram ti teser om typen, van de Velde svara med ti motteser om kunstnaren, og saka vart ikkje avgjord ved eit argument den eine måtte bøye seg for. Ho vart røysta over. Eit fag som må røyste over forma si, har ikkje teke feil av kva slag fag det er. Det har innrømt kva slag fag det er.

\section{«Den tause kunnskapen er kollektiv, ikkje privat»}

Den neste innvendinga rettar seg mot kapitlet om atelieret. Du hevdar at den tause kunnskapen i design er privat heile vegen ned, seier ho, men det er empirisk feil. Taus kunnskap er for det meste \emph{kollektiv}, tileigna gjennom sosialisering i ein praksisfellesskap \autocite{collins2010tacit}; og all kunnskap, fysikkens med, kviler på ein taus dimensjon \autocite{polanyi1966tacit}. Om ein taus kjerne gjorde eit fag grunnlaust, fall fysikken med design.

Innvendinga har rett om at taus kunnskap kan vere delt, og eg trekkjer gjerne «privat heile vegen ned» tilbake om det blir lese for sterkt. Men knowing how er ikkje same sak som knowing that, og skiljet er heile svaret. Fysikaren sin tause kunnskap er den private realiseringa av eit eksplisitt fundament: han har lagt i hendene noko som finst skrive ned, delt og prøvd utanfor hovudet hans. Kirurgen likeins, der den tause ferdigheita kviler på ein anatomi som ikkje er taus. Designaren har den tause realiseringa utan det eksplisitte fundamentet under. Og at noko er kollektivt taust, gjer det ikkje til kunnskap i den meininga eit fag treng. Mote er kollektivt taus. Ein heil generasjon kan dele ein smak utan at nokon kan grunngje han, og delinga gjer ikkje smaken til kunnskap; ho gjer han berre utbreidd. Det er nett dette atelieret overfører, ein delt smak som er kollektiv utan å vere prøvbar. Så innvendinga vinn det ho ber om: tausleiken treng ikkje vere einsleg. Berre at delt tausleik framleis kan vere smak, og at design ikkje har vist han er noko meir.

\section{«Det finst kumulativ kunnskap»}

Tredje innvendinga er den mest konkrete, og difor den eg tek tyngst. Du seier design ikkje akkumulerer, men sjå på moteksempla. Mønsterspråk som vandrar heilt frå arkitektur til programvare og ber \autocite{gof1994patterns}. Heile, replikerbare forskingsmetodologiar \autocite{blessing2009drm}. Empirisk testa, replikerte designpåstandar, som at naturutsyn frå sjukesenga kortar liggjetida \autocite{ulrich1984view}. Eit nordisk miljø som har vist korleis design lagar si eiga teori gjennom konstruktiv forsking og program og eksperiment \autocite{redstrom2017making, sevaldson2010discussions, koskinen2011constructive}, og som har namngjeve eit eige register mellom det einskilde artefaktet og teorien \autocite{lowgren2013intermediate}.

Dette kunne felt boka, om det heldt heilt ut. Det held ikkje, men det tvingar fram ein presisering eg gjev med glede. Sjå nøye på kva dei beste moteksempla faktisk er. Ulrich-studien er ekte, kumulativ kunnskap, men metoden og falsifiseringa og kumulasjonen er miljøpsykologien sin, brukt på ein designkontekst. Mønsterspråket som vandra til programvare, vart kumulativt i informatikken, der det fanst eit fundament å feste det i, ikkje i arkitekturen det kom frå. Dei nordiske forsøka på eigen teori er dei mest seriøse, og dei fortener respekt, men «program og eksperiment» og «intermediate-level knowledge» er framleis rammeverk for å \emph{organisere} praksis, ikkje ein delt, prøvbar modell av kva forma er. Eg legg merke til at sjølve namnet på Löwgrens kategori, kunnskap mellom det einskilde verket og teorien, er ei ærleg innrømming av kvar feltet står: det har det fyrste og strekkjer seg mot det andre utan å nå fram. Eit mønster går att, og det er verdt å seie kva det \emph{ikkje} viser. Der design verkeleg akkumulerer, gjer det det ved å låne ein annan kunnskaps fundament eller migrere kunnskapen dit eit fundament alt finst. Det viser ikkje at design er ringare for det. Det viser at form ikkje er det slag ting eit fundament av den typen kan gripe. Akkumulasjonen samlar seg kring forma, aldri i henne.

\section{«Vrangskapen er ei ekte innsikt, ikkje eit alibi»}

Fjerde innvendinga forsvarar omgrepet om vrange problem. Rittel, som mynta det, brukte det ikkje for å sleppe teori; han bygde rigorøse argumentasjonsmetodar som svar \autocite{rittel1973}. Innsikta er ekte. Designproblem er strukturelt ulike tamme problem, og å krevje naturvitskapleg falsifiserbarheit av dei er sjølv ein kategorifeil.

Eg gjev innvendinga heile kjernen. Vrangskapen er ekte, og Rittel meinte ikkje noko alibi. Men det var aldri Rittel eg klaga på, det var bruken faget gjorde av han, og her stadfestar innvendinga klagemålet i staden for å felle det. For ho innrømmer at vrangskapen krev ein \emph{annan} type teori, ein teori om kva ein gjer under motstridande krav, og det er nett ein slik teori faget ikkje har bygt. Å seie «problema er vrange, difor treng vi ein teori om vrangskap» er å gje meg rett. Å seie «problema er vrange, difor treng vi ingen teori» er alibiet. Faget har sagt det andre og sitert folk som meinte det fyrste. Innvendinga reddar Rittel. Han trong aldri reddast.

\section{«Design har likevel ein etikk»}

Femte innvendinga gjeld det eg sjølv har kalla fråvêret av eit etisk organ. Design har rike etikkar, modellkodeksar, profesjonsstandardar, deklarasjonar signerte i nærvær av \textsc{fn}-organ, ein heil tradisjon frå First Things First og framover. Påstanden om at faget manglar ei etikk, er rett og slett feil.

Her gjev eg innvendinga mest, og det skadar ikkje boka. Det finst rike etikkar i design, og eg skulle vore meir presis. Men sjå kva som står att når presisjonen er på plass. Alle desse kodeksane deler eitt drag: nesten ingen kan handhevast over individet. Arkitektane sitt organ kan frata ein utøvar retten til å praktisere \autocite{riba2019code}. Legen som bryt eden, kan misse retten til å vere lege. Designfaget sitt organ kan i beste fall stryke eit frivillig medlemskap, og den som strauk er framleis designar dagen etter. Skiljet mellom å frata retten og å stryke eit medlemskap er heile skiljet mellom ein profesjon med ryggrad og ein moralsk tradisjon utan tvang. Design har det andre i overflod og det fyrste nesten ikkje. Ein etikk ingen kan tvingast til å følgje, og som ingen mister noko på å bryte, disiplinerer ikkje faget. Han pyntar det.

\section{«Design har si eiga teori»}

Sjette innvendinga er den mest lærde. Du teiknar eit fag utan teori, men det er rein vankunne. Det finst ein eigen designvitskap, C-K-teorien, som modellerer designresonnering som ei utviding mellom eit rom av konsept og eit rom av kunnskap \autocite{hatchuel2009}. Det finst Design Science Research, eit heilt paradigme med eigne prosesssteg \autocite{hevner2004dsr, peffers2007dsrm}. Og det finst research through design, ein eigen kunnskapsveg der det å lage artefaktet \emph{er} forskinga \autocite{frayling1993, zimmerman2007rtd}.

Ta dei to som gjer reelt arbeid. C-K-teorien voks fram hjå Armand Hatchuel og Benoît Weil ved ingeniørdesign-pensumet på Mines ParisTech, skissert i 1996 og konsolidert kring 2003. Konsept-rommet C rommar proposisjonar som er avgjerdslause, utan logisk status i kunnskapsrommet K, og designoperasjonen er ei utviding mellom dei to. Det er ein vakker formalisme, og eg meiner ordet. Men sjå presist kva han modellerer. Han modellerer korleis ein resonnerer mot det ukjende, ikkje kva som gjer ei ferdig form rett. Han fortel deg kva ein designhandling er, ikkje kvifor denne stolen er betre enn den. Det er det andre svaret som er verdt å merke seg, for C-K kjem nettopp ut av ingeniørdesignen, og kan difor ikkje førast som vitne om at design \emph{burde} likne ingeniørfaget. Han syner heller at to nærskylde praksisar kan dele ein modell av tenkinga og likevel skiljast på kva slag dom dei skal felle. Design Science Research er det andre tilfellet, og namnet seier «science», men sjå kvar grunnen ligg. Han ligg i informasjonssystem, der ein artefakt anten verkar eller ikkje verkar mot eit definert problem, og det er køyringa som dømer \autocite{hevner2004dsr}. DSR akkumulerer fordi han har henta dommen sin frå eit felt der ein ferdig ting kan prøvast i drift. Det er ein prøve design godt kan låne, men han kjem utanfrå og dømer noko anna enn det design elles seier er kjernen sin: ikkje om forma er sann eller vakker, berre om ho verkar. Begge rammeverka er ekte. Begge svarar på eit anna spørsmål enn det boka stiller.

\begin{kasus}{Gaver tek ordet frå sitt eige fag}
Det mest avslørande vitnemålet kjem ikkje frå ein kritikar, men frå den fremste forsvararen. William Gaver, sjølv ein hovudfigur i research through design, sette seg føre å svare på kva ein kan vente av feltet, og kom til ein konklusjon som feller meir enn han reddar \autocite{gaver2012what}. RtD-kunnskap, skriv han, er og bør vere «provisional, contingent and aspirational». Generaliserte teoriar har for han ei rolle avgrensa til «inspiration and annotation»; «it is the artifacts we create that are the definite facts of research through design». Han åtvarar uttrykkjeleg mot å vente stabile, kumulative teoriar av feltet, fordi krav om det ville drepe nett det som gjer det verdfullt. Legg merke til kven som talar. Dette er ikkje ein motstandar som nektar feltet eit fundament; det er den næraste vennen som seier at fundamentet korkje finst eller bør etterspørjast. Jonas Löwgren prøver å berge noko midt imellom med annotated portfolios og mellomnivå-kunnskap, meir generell enn det einskilde verket og mindre enn ein teori \autocite{lowgren2013intermediate}. Men «mellom artefakt og teori» er sjølve adressa: feltet bur der, og strekkjer seg mot teorien utan å flytte inn. Eit fag som gjennom sin beste talsmann ber deg om \emph{ikkje} å vente kumulativ teori av det, har sagt sanninga om sin eigen grunn. Eg gjer ikkje anna enn å sitere han.
\end{kasus}

\section{«Det finst evidensbasert design»}

Den siste innvendinga er den hardaste empiriske, og eg tek henne sist fordi ho ser ut til å felle alt med ein einaste solid kunnskapskropp. Du seier design ikkje har ein eigen, prøvbar, kumulativ vitskap, men evidensbasert design finst og er moge. I helsearkitektur er det vist og replikert at romutforming verkar på pasientutfall, og feltet har vakse til ein kunnskapskropp med systematiske oversyn over korleis det fysiske miljøet verkar på helse \autocite{ulrich2008ebd}. Her gjer design påstandar, prøver dei, og byggjer kumulativt.

Dette er det sterkaste moteksempelet i boka, og difor det viktigaste å få rett. Ta Ulrich-studien presist, for talet ber svaret. Pasientar med utsyn til natur frå senga hadde kortare liggjetid og bad om mindre sterke smertestillande enn dei med utsyn til ein murvegg. Kontrollgruppe, måleeiningar i liggjedøgn og dosar, ei falsifiserbar hypotese \autocite{ulrich1984view}. Spør kven som eig kvar del. Falsifiseringa er psykologien sin. Kontrollgruppa er medisinen si. Statistikken og dei systematiske oversyna er helseforskinga si kumulasjonsform. Det einaste design bidreg med, er gjenstanden som blir prøvd, vindauget. Metoden er lånt heilt, gjenstanden er designs eigen, og det er nett dette delet som er poenget. Der design verkeleg akkumulerer slik, gjer det det ved å la ein framand metode døme ein designgjenstand. Eg skal ikkje seie at design då «endeleg» har nådd noko, som om vitskapen var toppen av ein stige design klatra. Det ville vere å smugle inn att malen eg nettopp avviste. Eg seier berre kor lite dette dekkjer. Det verkar der utfallet er målbart i kropp og tal, i liggjetid og infeksjonsrate og fallfrekvens, og det teier overalt elles, om kva som gjer ei skrift vakker, ei merkevare sann, ein stol rett. Evidensbasert design er ikkje design som fann sitt grunnlag. Det er design som låner eit anna felts grunnlag på det smale stykket der gjenstanden gjev frå seg eit tal, og som ingenting har å låne der gjenstanden ikkje gjer det.

\vignett

Sju innvendingar, kvar i si sterkaste form. To av dei rispa boka, og det skal stå: kategorifeilen har rett i at design ikkje er ein vitskap som kom til kort, og Gaver minner meg om at feltets beste folk aldri lova det grunnlaget eg etterlyser. Eg kan ikkje bevise at form \emph{må} kunne grunngjevast slik. Det eg kan vise, gjekk uskadd gjennom alle sju. Kvar gong eit «fundament» vart peika på, synte det seg å vere anten ein modell av korleis ein tenkjer eller eit lån frå ein annan kunnskap, og aldri ein delt, prøvbar grunn til å seie kvifor denne forma er betre enn den. Det fråvêret gjer ikkje design til ein dårleg vitskap. Det gjer design til noko anna, eit fag som må røyste der andre kan måle, og som difor lyt grunngjevast på sine eigne premiss, ikkje på lånte. Boka si oppgåve var aldri å straffe faget for det. Ho var å seie det høgt.

\laeringsmal{\begin{itemize}
\item Attgje dei sterkaste forsvara for designfaget i si beste form, frå kategorifeil til evidensbasert design.
\item Skilje ein prosessformalisme, ein teori om korleis ein tenkjer, frå eit grunnlag for å døme den ferdige forma.
\item Forklare kvifor «design er ikkje fysikk» er sant og likevel ikkje svarar på kniven boka stiller.
\end{itemize}}

\ovingar{\begin{enumerate}
\item Vel den innvendinga du finn sterkast. Før henne i si beste form, og prøv så å felle bokas svar.
\item Kategorifeil-innvendinga og Gaver-vitnemålet fekk lov å rispe boka. Vel ei av dei og avgjer: rispar ho, eller feller ho? Grunngjev.
\item Konstruer ein åttande innvending boka ikkje tek opp. Held han mot kniven, eller går han, som dei andre, utanom henne?
\end{enumerate}}

\vidarelesing{\fullcite{cross2001discipline}; \fullcite{collins2010tacit}; \fullcite{gof1994patterns}; \fullcite{ulrich1984view}; \fullcite{rittel1973}; \fullcite{gaver2012what}; \fullcite{riba2019code}.}
